\section{Invertibility of categories of sheaves}
\label{sec:inv_sheaves}

We now ask when categories of sheaves are invertible in the higher Morita categories of Section
\ref{sec:Morita}. For the symmetric monoidal structure the answer is essentially never (Section
\ref{sec:inv_QC}). We therefore turn to convolution structures (Section \ref{sec:convolution}),
and in particular to quasi-coherent sheaves on loop stacks, which are Morita trivial whenever the
stack satisfies diagonal descent (Section \ref{sec:diag_descent}). This includes the formal
completion of a vector space at the origin (Section \ref{sec:formal_vector_space}).

\subsection{Symmetric monoidal categories of sheaves}
\label{sec:inv_QC}

We work relative to a base stack $S$; the absolute case is $S=\pt$. For a map $f\colon X\to S$,
the unit of the commutative $\QC\left(S\right)$-algebra $\left(\QC\left(X\right),\otimes\right)$ is
$f^*\colon\QC\left(S\right)\to\QC\left(X\right)$, so by Corollary \ref{cor:inv_symm_rel} it is
invertible in $\umor_d\left(\Mod_{\QC\left(S\right)}\left(\Pre\right)\right)$ if and only if $f^*$ is
an equivalence. To translate this into a condition on stacks, we use Tannaka duality.
\begin{lem}
\label{lem:Tannaka}
Let $f\colon X\to Y$ be a map of stacks such that $f^*$ is an equivalence.
\begin{enumerate}
\item If $f$ has a retraction $g$, i.e.\ $gf\simeq\Id_X$, and $Y$ is quasi-compact with
affine diagonal, then $f$ is an equivalence.
\item If $f$ has a section $s$, i.e.\ $fs\simeq\Id_Y$, and $X$ is quasi-compact with
affine diagonal, then $f$ is an equivalence.
\item If $f$ is flat and $X$ and $Y$ are quasi-compact with affine diagonal, then $f$ is an
equivalence.
\end{enumerate}
\end{lem}

\begin{proof}
For (1): $f^*g^*\simeq\Id$, so $g^*\simeq\left(f^*\right)^{-1}$, and hence
$\left(fg\right)^*\simeq g^*f^*\simeq\Id$ as symmetric monoidal functors.
By Tannaka duality \cite[Theorem 9.2.0.2]{L:SAG}, \cite{BHL}, maps into $Y$
are detected by their pullback functors, so $fg\simeq\Id_Y$.
Part (2) is the same argument with the roles reversed: $s^*\simeq\left(f^*\right)^{-1}$,
so $\left(sf\right)^*\simeq\Id$ and $sf\simeq\Id_X$.

For (3): Since $f$ is flat, $f^*$ is t-exact, and being conservative it detects connective
and discrete objects: If $f^*E$ is connective, then
$f^*\tau_{<0}E\simeq\tau_{<0}f^*E\simeq 0$, so $\tau_{<0}E\simeq 0$, and similarly for
$\tau_{>0}$.
Let $F\ceqq\left(f^*\right)^{-1}$. Then $F$ preserves connective objects, since
$f^*F\simeq\Id$. It also preserves flat objects, using that a connective object $E$ is flat if and only if
$E\otimes M$ is discrete for every discrete $M$ (affine-locally, \cite[Section 7.2.2]{L:HA}): If $E$ is flat
and $M\in\QC\left(Y\right)$ is discrete, then $f^*\left(F\left(E\right)\otimes M\right)\simeq
E\otimes f^*M$ is discrete, hence so is $F\left(E\right)\otimes M$.
So $F$ is a colimit-preserving symmetric monoidal functor preserving connective and flat
objects, and by Tannaka duality it is $h^*$ for some $h\colon Y\to X$.
Then $\left(fh\right)^*\simeq h^*f^*\simeq\Id$ and $\left(hf\right)^*\simeq f^*h^*\simeq\Id$,
so $fh\simeq\Id_Y$ and $hf\simeq\Id_X$.
\end{proof}

\begin{theorem}
\label{thm:inv_rel}
Let $f\colon X\to S$ be a map of quasi-compact stacks with affine diagonal, and let
$d\geq 1$.
Suppose that $f$ is flat, or that $f$ has a section, or that $f$ has a retraction.
Then the commutative $\QC\left(S\right)$-algebra $\left(\QC\left(X\right),\otimes\right)$
is invertible in $\umor_d\left(\Mod_{\QC\left(S\right)}\left(\Pre\right)\right)$ if and only if
$f$ is an equivalence.
\end{theorem}

\begin{proof}
By Corollary \ref{cor:inv_symm_rel}, invertibility is equivalent to $f^*$ being an
equivalence, and by Lemma \ref{lem:Tannaka} this is equivalent to $f$ being an equivalence.
\end{proof}

\subsection{Convolution categories and descent}
\label{sec:convolution}
\label{sec:conv}

When the stack on which we study sheaves is a fiber product, convolution gives
noncommutative monoidal structures.

Given a map of stacks $X\to Y$, the category of sheaves on the fiber product $X\times_Y X$
has a monoidal structure given by convolution.
Namely, we have a diagram of stacks:
\begin{equation*}
\begin{tikzcd}[column sep=small]
 & X\times_Y X\times_Y X
   \arrow[dl, "p_{12}"']
   \arrow[d, "p_{13}"]
   \arrow[dr, "p_{23}"] & \\
 X\times_Y X & X\times_Y X & X\times_Y X
\end{tikzcd}
\end{equation*}
and convolution
is induced by pulling and pushing along this diagram:
\begin{equation*}
M* N = p_{13*}\left(p_{12}^*M\otimes p_{23}^*N\right)
\end{equation*}

\begin{exm}
When $Y = \pt$, $\left(\QC\left(X\times X\right),*\right)$ is the monoidal category of integral
kernels, equivalent for $X$ perfect to $\End\left(\QC\left(X\right)\right)$ under composition
(Remark \ref{rmk:BZFN_monoidal}).
\end{exm}

\begin{theorem}[{\cite[Theorem 1.3]{BZFN:app}}]
Let $\pi \colon X \to Y$ be a map of perfect stacks.
If either
\begin{itemize}
\item $\pi$ is faithfully flat, or
\item $X$ and $Y$ are smooth and $\pi$ is proper and surjective,
\end{itemize}
then tensoring with
the $\QC\left(Y\right)$-$\QC\left(X\times_Y X\right)$-bimodule $\QC\left(X\right)$ defines
an equivalence
\begin{equation*}
\Mod_{\left(\QC\left(Y\right) , \otimes\right)}\left(\Pre\right) \lto{\sim}
\Mod_{\left(\QC\left(X\times_Y X\right) , *\right)}\left(\Pre\right)
\end{equation*}
\label{thm:BZFN_app}
\end{theorem}

\subsection{Quasi-coherent sheaves on loop stacks}

A particularly important instance of a convolutional category of sheaves on a fiber
product is given by the diagonal map $X \to Y = X\times X$.
This is the category obtained via Theorem \ref{thm:BZFN_loops} as the Drinfeld center of
the symmetric-monoidal category $\QC\left(X\right)$.

\subsubsection{Descent along the diagonal}
\label{sec:diag_descent}

\begin{defn}
\label{defn:diag_descent}
A stack $X$ satisfies \emph{diagonal descent} if
\begin{enumerate}
\item the conclusions of Theorem \ref{thm:BZFN} hold for the diagonal: $\QC\left(X\right)$ is
self-dual in $\Pre$ and over $\QC\left(X^2\right)$,
$\QC\left(X\right)\otimes\QC\left(X\right)\simeq\QC\left(X^2\right)$, and
\begin{equation*}
\QC\left(X\right)\otimes_{\QC\left(X^2\right)}\QC\left(X\right)\simeq\QC\left(\cL X\right)
\simeq\End_{\QC\left(X^2\right)}\left(\QC\left(X\right)\right) \ ,
\end{equation*}
where the second equivalence identifies composition with convolution; and
\item the conclusion of Theorem \ref{thm:BZFN_app} holds for the diagonal, i.e.\ tensoring
with $\QC\left(X\right)$ gives an equivalence
\begin{equation*}
\Mod_{\left(\QC\left(X^2\right),\otimes\right)}\left(\Pre\right)\simeq
\Mod_{\left(\QC\left(\cL X\right),*\right)}\left(\Pre\right) \ .
\end{equation*}
\end{enumerate}
\end{defn}

Condition (1) holds for every perfect stack by Theorem \ref{thm:BZFN} and Remark
\ref{rmk:BZFN_monoidal}, and it implies $\mathcal{Z}_1\left(\QC\left(X\right)\right)\simeq\QC\left(\cL X\right)$
as in Theorem \ref{thm:BZFN_loops}.
By Theorem \ref{thm:BZFN_app}, a perfect stack satisfies (2), and hence diagonal descent,
if $\Delta\colon X\to X\times X$ is faithfully flat, or if $X$ is smooth and $\Delta$ is
proper and surjective.

\begin{exm}
\label{exm:BG_diag_descent}
If $G$ is an affine algebraic group of finite type over a field of characteristic zero, then
$BG$ is perfect \cite[\S 3.3]{BZFN} and satisfies diagonal descent.
Indeed, the base change of $\Delta$ along the cover $\pt\to BG\times BG$ is
$\left(G\times G\right)/G\simeq G\to\pt$, so $\Delta$ is faithfully flat.
\end{exm}

\begin{rmk}
\label{rmk:diag_descent_fails}
Diagonal descent fails for every quasi-compact separated scheme $X$ with more than one point,
e.g.\ any positive-dimensional scheme of finite type over $k$.
Indeed, $X$ is perfect, and its diagonal is not surjective, so there is a non-empty affine open
$U\subset X^2$ disjoint from $\Delta\left(X\right)$. By Theorem \ref{thm:BZFN},
\begin{equation*}
\QC\left(X\right)\otimes_{\QC\left(X^2\right)}\QC\left(U\right)\simeq
\QC\left(X\times_{X^2}U\right)\simeq\QC\left(\emptyset\right)=0 \ ,
\end{equation*}
whereas $\QC\left(U\right)\neq0$. So tensoring with $\QC\left(X\right)$ is not conservative,
and condition (2) of Definition \ref{defn:diag_descent} fails.
\end{rmk}

The formal completion of a vector space at the origin also satisfies diagonal descent,
although it is not perfect and its diagonal is not flat (Proposition
\ref{prop:Morita_formal}).
Stacks satisfying diagonal descent are exactly those to which the Morita triviality result,
Theorem \ref{thm:trivial_QC}, applies.

\subsubsection{Invertibility}

\begin{lem}[Morita context]
\label{lem:Morita_context}
Let $X$ be a stack satisfying condition (1) of Definition \ref{defn:diag_descent}, and write
$R\ceqq\QC\left(X^2\right)$, $M\ceqq\QC\left(X\right)$ and
$A\ceqq\left(\QC\left(\cL X\right),*\right)\simeq\End_R\left(M\right)$. Then
$\Psi\ceqq M\otimes_A-\colon\Mod_A\left(\Pre\right)\to\Mod_R\left(\Pre\right)$ is left adjoint
to $\Phi\ceqq M^\vee\otimes_R-$ and fully faithful, and
\begin{equation*}
P\ceqq M\otimes_A M^\vee=\QC\left(X\right)\otimes_{\QC\left(\cL X\right)}\QC\left(X\right)
\lto{\ \varepsilon\ } R
\end{equation*}
is an idempotent $R$-coalgebra, i.e.\ $P\otimes_R P\simeq P$ via $\varepsilon$. The functor $\Phi$ is an equivalence if and only if the counit
$\varepsilon$ is an equivalence, which is condition (2) of Definition \ref{defn:diag_descent}.
\end{lem}

\begin{proof}
Since $M$ is dualizable over $R$ with dual $M^\vee\simeq M$, the right adjoint of $\Psi$ is
$\Fun_R\left(M,-\right)\simeq\Phi$, and $M^\vee\otimes_R M\simeq A$ as $A$-bimodules, so the
unit $\Id\to\Phi\Psi$ is an equivalence. In the other order, $\Psi\Phi\simeq P\otimes_R-$, with
counit $\varepsilon\otimes_R-$, and
$P\otimes_R P\simeq M\otimes_A\left(M^\vee\otimes_R M\right)\otimes_A M^\vee\simeq P$.
\end{proof}

The following upgrades the Morita equivalence of \cite{BZFN:app} in $\umor_1\left(\Pre\right)$
(Theorem \ref{thm:BZFN_app}) to $\umor_2\left(\Pre\right)$, under the different hypothesis of diagonal
descent, using \cite{BZFN} and the criterion of \cite[Theorem 2.26]{BJSS}.

\begin{theorem}
\label{thm:trivial_QC}
Let $X$ be a stack satisfying diagonal descent (Definition \ref{defn:diag_descent}), e.g.\
$X = BG$ (Example \ref{exm:BG_diag_descent}) or the formal completion $\widehat V_0$ of a
vector space at the origin (Proposition \ref{prop:Morita_formal}).
Then $\left(\QC\left(\mathcal{L}X\right),*\right)$ is Morita trivial, i.e.\ equivalent to the
unit of $\umor_2\left(\Pre\right)$.
More precisely, $\QC\left(X\right)$, regarded as a $\QC\left(\cL X\right)$-central algebra via
$\mathcal{Z}_1\left(\QC\left(X\right)\right)\simeq\QC\left(\cL X\right)$, is an invertible $1$-morphism
from $\QC\left(\mathcal{L}X\right)$ to the monoidal unit of $\umor_2\left(\Pre\right)$.
\end{theorem}

\begin{proof}
We check the criterion of Theorem \ref{thm:trivial_Alg2} \cite[Theorem 2.26]{BJSS}, for
$B = \left(\QC\left(\cL X\right),*\right)$ and $C = \left(\QC\left(X\right),\otimes\right)$. Note
that $C^e = C\otimes C^\rev\simeq\left(\QC\left(X^2\right),\otimes\right)$ since $C$ is
commutative.
By condition (1) of Definition \ref{defn:diag_descent}, $C$ is dualizable in $\Pre$ and over
$C^e$, i.e.\ $2$-dualizable in $\umor_1\left(\Pre\right)$, and $C$ is self-dual in both senses,
so $C^!\simeq C^\vee\simeq\QC\left(X\right)$.
\begin{enumerate}
\item The evaluation $C^!\otimes_B C\to C^e$ is the counit
$\QC\left(X\right)\otimes_{\QC\left(\cL X\right)}\QC\left(X\right)\to\QC\left(X^2\right)$ of
Lemma \ref{lem:Morita_context}. It is an equivalence
by condition (2).
\item The map $B\to \mathcal{Z}_1\left(C\right)$ is the equivalence of condition (1), cf.\ Theorem
\ref{thm:BZFN_loops}.
\item[(4)] Since $C$ is commutative, $C^\rev\otimes_B C\to\End\left(C\right)$ is the same counit,
followed by the equivalence $\QC\left(X^2\right)\simeq\Fun\left(\QC\left(X\right),\QC\left(X\right)\right)$
given by self-duality of $\QC\left(X\right)$ in $\Pre$. So it is an equivalence.
\item[(3)] By (4), $C\otimes_B C^\rev\simeq\End\left(C\right)$, and the evaluation in (3)
becomes $C^\vee\otimes_{\End\left(C\right)}C\to\Vect$. This is an equivalence for any dualizable
$C\in\Pre$, since $\End\left(C\right)\simeq C\otimes C^\vee$ is Morita trivialized by $C$. \qedhere
\end{enumerate}
\end{proof}

In particular, by Theorem \ref{thm:inv_E2}, $\left(\QC\left(\cL X\right),*\right)$ is
non-degenerate, factorizable and cofactorizable. For instance, for $X=BG$, where $\cL BG\simeq G/G$ is the adjoint quotient, we have
\begin{equation*}
\mathcal{Z}_2\left(\QC\left(G/G\right),*\right)\simeq\Vect
\qquad\text{and}\qquad
\QC\left(G/G\times G/G\right)\simeq \mathcal{Z}_1\left(\QC\left(G/G\right),*\right) \ ,
\end{equation*}
the latter via the action map, generalizing the modularity of the Drinfeld center of
$\Rep\left(G\right)$ for finite $G$.

\subsubsection{The formal completion of a vector space at the origin}
\label{sec:formal_vector_space}

Let $V$ be a finite-dimensional vector space and $\widehat V_0\ceqq V^\wedge_0$ its formal
completion at the origin. Strictly speaking $\widehat V_0$ is an ind-scheme rather than a
perfect stack, so Theorems \ref{thm:BZFN} and \ref{thm:BZFN_app} do not literally apply to
it. Nonetheless it satisfies diagonal descent.

\begin{prop}
\label{prop:Morita_formal}
The formal completion $\widehat V_0$ satisfies diagonal descent
(Definition \ref{defn:diag_descent}). In particular, tensoring with $\QC\left(\widehat V_0\right)$
gives a Morita equivalence between $\left(\QC\left(\widehat V_0^2\right),\otimes\right)$ and
$\left(\QC\left(\cL\widehat V_0\right),*\right)$, and
$\mathcal{Z}_1\left(\QC\left(\widehat V_0\right)\right)\simeq\QC\left(\cL\widehat V_0\right)$.
\end{prop}

\begin{proof}
\emph{Step 0: Completions as supports.}
For a derived scheme $Y$ of finite type and a closed point $0\in Y$, local duality identifies
$\QC\left(Y^\wedge_0\right)$, symmetric monoidally, with the category $\QC_0\left(Y\right)$ of
quasi-coherent sheaves set-theoretically supported at $0$
\cite[Chapters 7--8]{L:SAG}. The local cohomology functor $\Gamma_0$ is
smashing, $\Gamma_0 F\simeq F\otimes\Gamma_0\cO_Y$, so $\QC_0\left(Y\right)$ is an idempotent
commutative $\QC\left(Y\right)$-algebra, and for $f\colon Z\to Y$ we have
$\QC\left(Z\right)\otimes_{\QC\left(Y\right)}\QC_0\left(Y\right)\simeq\QC_{f^{-1}\left(0\right)}\left(Z\right)$.
Completion commutes with fiber products, so $\widehat V_0^n\simeq\left(V^n\right)^\wedge_0$ and
$\cL\widehat V_0\simeq\left(\cL V\right)^\wedge_0=\widehat V_0\times V\left[-1\right]$.
Hence $\QC\left(\widehat V_0^n\right)\simeq\QC_0\left(V^n\right)$ and
$\QC\left(\cL\widehat V_0\right)\simeq\QC_0\left(\cL V\right)$.

\emph{Step 1: Condition (1).}
Write $E\ceqq\QC_0\left(V^2\right)$, $M\ceqq\QC_0\left(V\right)$ and
$A\ceqq\QC_0\left(\cL V\right)$. Since the origin of $V^2$ has preimage the origin under both
$\Delta\colon V\to V^2$ and $\cL V\to V^2$, Step 0 gives
$M\simeq\QC\left(V\right)\otimes_{\QC\left(V^2\right)}E$ and
$A\simeq\QC\left(\cL V\right)\otimes_{\QC\left(V^2\right)}E$; moreover $M\otimes M\simeq E$.
Base change along the symmetric monoidal functor $\QC\left(V^2\right)\to E$ preserves duals,
relative tensor products, and endomorphisms of dualizable modules. So Theorem
\ref{thm:BZFN} for $\Delta\colon V\to V^2$, together with Remark \ref{rmk:BZFN_monoidal},
gives condition (1) for $\widehat V_0$: $M$ is self-dual over $E$, and
$M\otimes_E M\simeq A\simeq\End_E\left(M\right)$. Finally, $M$ is compactly generated by the
Koszul complex $K_V$ of the origin, so $M\simeq\Ind\left(\Perf_0\left(V\right)\right)$, where
$\Perf_0\left(V\right)$ denotes perfect complexes set-theoretically supported at the origin.
Since $\mathcal{H}om\left(-,\cO_V\right)$ preserves $\Perf_0\left(V\right)$, we have
$\Perf_0\left(V\right)^{\op}\simeq\Perf_0\left(V\right)$ and $M$ is self-dual in $\Pre$.

\emph{Step 2: Condition (2).}
In Lemma \ref{lem:Morita_context}, $\Psi = M\otimes_A-$ is left adjoint to
$\Phi = M^\vee\otimes_E- = \Hom_E\left(M,-\right)$, and it is fully faithful. So the counit
$\varepsilon\colon P=\Psi\Phi\left(E\right)\to E$ is an equivalence as soon as $E$ lies in the
essential image of $\Psi$, which we now show.
Let $\Delta^\times\colon E\to M$ be the right adjoint of $\Delta_*\colon M\to E$. Since $\Delta$ is a
regular closed immersion of codimension $\dim V$, we have
$\Delta^\times\simeq\Delta^*\left(-\right)\otimes\det\left(N_\Delta\right)\left[-\dim V\right]$, so $\Delta^\times$
is continuous and $E$-linear.
It is also conservative, and this is where completing matters: $\Delta^\times$ is not
conservative on $\QC\left(V^2\right)$, but on $E$ it is. Indeed, the compact generator
$K_{V^2}$ of $E$ is $\Delta_*K_V$, since the ideals $\left(x_i,y_i\right)$ and
$\left(x_i-y_i,x_i\right)$ agree. So if $\Delta^\times F = 0$ then
$\Hom\left(K_{V^2},F\right)\simeq\Hom\left(K_V,\Delta^\times F\right)=0$, and $F=0$.
By the Barr--Beck--Lurie theorem \cite[Theorem 4.7.3.5]{L:HA}, $\Delta^\times$ identifies $E$ with
the category $\Mod_T\left(M\right)$ of modules over the monad $T=\Delta^\times\Delta_*$. This monad
is $E$-linear, so it is an algebra in $\End_E\left(M\right)\simeq A$, and
$\Mod_T\left(M\right)\simeq M\otimes_A\Mod_T\left(A\right)$ \cite[Theorem 4.8.4.6]{L:HA}.
Hence $E\simeq\Psi\left(\Mod_T\left(A\right)\right)$ lies in the essential image of $\Psi$.
\end{proof}
